\documentclass[UTF8]{ctexart}
\usepackage{amsmath}
\usepackage{amssymb}
\usepackage{geometry}
\usepackage{enumerate}
\geometry{a4paper, margin=1in}

\title{计算机图形学计算专题练习（110题）}
\author{整理：Antigravity}
\date{\today}

\begin{document}

\maketitle

\section{DDA 直线生成算法}
\begin{enumerate}
    \item 计算从 $(0,0)$ 到 $(8,3)$ 的直线。写出每步的 $\Delta x, \Delta y$ 及四舍五入后的像素坐标。
    \item 计算从 $(2,1)$ 到 $(10,5)$ 的直线。
    \item 计算从 $(-2,-1)$ 到 $(3,2)$ 的直线。
    \item 计算从 $(0,0)$ 到 $(3,8)$ 的直线（注意主位移方向）。
    \item 计算从 $(5,5)$ 到 $(2,2)$ 的直线。
    \item 计算从 $(0,0)$ 到 $(-5,2)$ 的直线。
    \item 计算从 $(1,2)$ 到 $(6,1)$ 的直线。
    \item 计算从 $(0,0)$ 到 $(4,4)$ 的直线。
    \item 计算从 $(10,10)$ 到 $(15,12)$ 的直线。
    \item DDA 算法在绘制斜率 $k > 1$ 的直线时，为什么要交换步进的方向？请以 $(0,0)$ 到 $(2,5)$ 为例说明。
\end{enumerate}

\section{中点画线法}
\begin{enumerate}
    \item 绘制从 $(0,0)$ 到 $(5,2)$ 的直线，写出初始判别式 $d_0$ 及每步的选择。
    \item 绘制从 $(1,1)$ 到 $(6,4)$ 的直线。
    \item 绘制从 $(0,0)$ 到 $(7,3)$ 的直线。
    \item 绘制从 $(2,2)$ 到 $(8,5)$ 的直线。
    \item 若直线斜率为 $0.5$，起点为 $(0,0)$，计算前 5 个点。
    \item 推导中点画线法中 $0 < k < 1$ 时误差项 $d$ 的递归公式。
    \item 说明中点画线法如何消除浮点运算。
    \item 计算从 $(-1,-1)$ 到 $(4,1)$ 的直线。
    \item 绘制从 $(0,0)$ 到 $(5,5)$ 的直线。
    \item 绘制从 $(10,5)$ 到 $(15,7)$ 的直线。
\end{enumerate}

\section{Bresenham 直线生成算法}
\begin{enumerate}
    \item 使用 Bresenham 算法绘制从 $(0,0)$ 到 $(5,2)$ 的直线，要求显示误差项 $d$ 变化过程。
    \item 绘制从 $(0,0)$ 到 $(2,5)$ 的直线（斜率 $k > 1$ 情况）。
    \item 绘制从 $(0,0)$ 到 $(-5,2)$ 的直线（第二象限）。
    \item 绘制从 $(0,0)$ 到 $(5,-2)$ 的直线（第四象限）。
    \item 绘制从 $(0,0)$ 到 $(-2,-5)$ 的直线（第三象限，$|k| > 1$）。
    \item 计算从 $(1,1)$ 到 $(6,3)$ 的直线。
    \item 计算从 $(2,1)$ 到 $(7,5)$ 的直线。
    \item 简述 Bresenham 算法相较于 DDA 算法的优点。
    \item Bresenham 算法在处理全象限直线时，需要考虑哪些变换？
    \item 计算从 $(-3,2)$ 到 $(2,4)$ 的直线。
\end{enumerate}

\section{中点画圆法}
\begin{enumerate}
    \item 半径 $R=5$，计算第一八分圆（$0 \le x \le y$）的像素点。
    \item 半径 $R=10$，计算第一八分圆的像素点。
    \item 半径 $R=8$，计算第一八分圆。
    \item 若圆心为 $(2,3)$，半径 $R=5$，计算其在第一八分圆对应的像素点。
    \item 说明圆的八对称性，并给出点 $(x,y)$ 对应的其余 7 个对称点坐标。
    \item 半径 $R=6$，写出初始判别式 $d_0$ 及第一步更新。
    \item 计算半径 $R=7$ 时，$x$ 从 $0$ 增加到 $x \approx y$ 的过程。
    \item 为什么中点画圆法只需计算 $45^\circ$ 的弧段？
    \item 半径 $R=12$，计算前 4 个点。
    \item 推导中点画圆法中 $d$ 的更新增量。
\end{enumerate}

\section{Bresenham 画圆算法}
\begin{enumerate}
    \item 半径 $R=5$，计算第一象限的像素点，写出误差项 $D$。
    \item 半径 $R=4$，写出完整的像素序列（利用对称性）。
    \item 半径 $R=10$，计算前 5 个像素点。
    \item 描述 Bresenham 画圆中，当点落在圆外、圆内和圆周上时，下一步的走向选择。
    \item 比较中点画圆法与 Bresenham 画圆法的异同。
    \item 半径 $R=6$，起始点为 $(0,6)$，计算下一步。
    \item 半径 $R=8$，计算第一象限。
    \item 半径 $R=3$，计算全圆路径。
    \item 若 $R=5$，写出当误差项 $d < 0$ 时的 $d$ 值更新公式。
    \item 计算半径 $R=9$，第一象限的像素点。
\end{enumerate}

\section{扫描线填充算法}
\begin{enumerate}
    \item 已知多边形顶点为 $(2,2), (8,2), (8,8), (2,8)$。构造 $y=4$ 时的 AET 表。
    \item 构造三角形 $(1,1), (5,1), (3,5)$ 的完整的 ET 表。
    \item 在扫描线填充中，如何处理多边形的水平边？
    \item 已知四边形 $(0,0), (4,0), (6,4), (2,4)$。计算 $y=2$ 时的交点并排序。
    \item 说明“下闭上开”原则在扫描线填充中的作用。
    \item 构造凹多边形 $(1,1), (5,1), (5,5), (3,3), (1,5)$ 的 ET 表。
    \item 已知扫描线 $y$ 与边相交，若交点坐标为浮点数，如何取整？
    \item 比较扫描线填充与种子填充算法的优缺点。
    \item 在 $y=10$ 时，AET 中有四条边，其交点 $x$ 坐标分别为 $2.5, 8.2, 10.1, 15.7$。写出填充的区间。
    \item 举例说明什么是“奇增偶减”原则。
\end{enumerate}

\section{2D 几何变换}
\begin{enumerate}
    \item 将点 $(1,1)$ 平移 $(2,3)$，写出齐次坐标下的变换矩阵和结果。
    \item 将点 $(2,2)$ 绕原点逆时针旋转 $90^\circ$。
    \item 将线段 $A(1,1), B(3,1)$ 以原点为中心缩放 2 倍。
    \item \textbf{复合变换}：将点 $(1,1)$ 绕点 $(2,2)$ 顺时针旋转 $90^\circ$。写出分解步骤。
    \item 对点 $(1,2)$ 进行沿 $x$ 轴方向的错切（$sh_x = 2$）。
    \item 求点 $(3,4)$ 关于直线 $y=x$ 的对称点矩阵。
    \item 证明两次连续的平移变换满足交换律。
    \item 证明旋转变换与缩放变换通常不满足交换律（给出反例）。
    \item 将三角形 $(0,0), (2,0), (1,2)$ 关于坐标原点进行镜像变换。
    \item 写出将窗口 $(x_{wmin}, y_{wmin}, x_{wmax}, y_{wmax})$ 映射到视口 $(x_{vmin}, y_{vmin}, x_{vmax}, y_{vmax})$ 的变换矩阵。
\end{enumerate}

\section{3D 几何变换}
\begin{enumerate}
    \item 写出绕 $Z$ 轴旋转 $\theta$ 角的齐次坐标变换矩阵。
    \item 将点 $(1,1,1)$ 平移 $(1,2,3)$。
    \item 将点 $(1,0,0)$ 绕 $Y$ 轴旋转 $90^\circ$。
    \item 计算绕直线 $x=1, y=1$（平行于 $Z$ 轴）旋转 $180^\circ$ 的变换矩阵。
    \item 写出绕 $X$ 轴旋转后，点 $(0,1,0)$ 变为 $(0,0,1)$ 的旋转角度。
    \item 描述三维坐标系下，右手系与左手系在旋转方向上的差异。
    \item 写出 $Z$ 轴方向的缩放矩阵（比例因子为 $s_z$）。
    \item 将点 $(x,y,z)$ 关于 $xoy$ 面进行镜像变换。
    \item 什么是透视投影矩阵？写出其最简形式。
    \item 复合变换：先绕 $X$ 轴旋转 $90^\circ$，再平移 $(1,1,1)$，最后绕 $Z$ 轴旋转 $90^\circ$。写出复合矩阵表达式（列向量表示）。
\end{enumerate}

\section{Cohen-Sutherland 裁剪算法}
\begin{enumerate}
    \item 裁剪窗口为 $[0,10] \times [0,10]$。点 $P( -1, 5)$ 的编码是什么？
    \item 若线段两端点编码分别为 $0001$ 和 $1000$，该线段是“简取”、“简弃”还是需进一步判断？
    \item 窗口为 $[0,5] \times [0,5]$，线段 $P_1(-1,2), P_2(6,2)$。计算与左边界的交点。
    \item 详细描述 Cohen-Sutherland 算法中“与”运算为非零时的含义。
    \item 给定 5 个点，分别位于窗口内、上、左、右下、左上，写出它们的区域编码。
    \item 线段 $P_1(1,1), P_2(8,8)$，窗口为 $[0,5] \times [0,5]$。写出裁剪过程。
    \item 为什么编码顺序（上、下、右、左）在算法中很重要？
    \item 在求交点时，若线段垂直于坐标轴，如何处理？
    \item 分析该算法在处理大型多边形裁剪时的局限性。
    \item 窗口为 $[0,10] \times [0,10]$，线段 $P_1(12,12), P_2(15,8)$。判断其位置关系。
\end{enumerate}

\section{Liang-Barsky 裁剪算法}
\begin{enumerate}
    \item 简述 $p_k < 0, p_k > 0$ 和 $p_k = 0$ 在算法中的物理意义。
    \item 已知线段 $P_1(1,1), P_2(4,3)$，窗口 $[0,2] \times [0,2]$。计算 $p_k$ 和 $q_k$。
    \item 在 Liang-Barsky 算法中，如何根据 $u_1, u_2$ 的大小判断线段是否在窗口外？
    \item 计算平行于 $x$ 轴的线段在 Liang-Barsky 算法中的处理流程。
    \item 线段 $P_1(-1, 5), P_2(3, -2)$，窗口 $[0,2] \times [0,2]$。计算进入参数 $u_1$。
    \item 比较 Liang-Barsky 与 Cohen-Sutherland 算法在效率上的差异。
    \item 推导 Liang-Barsky 算法中参数 $u$ 的求算公式。
    \item 当线段一个端点在窗口内，另一个在窗口外时，$u_1$ 和 $u_2$ 的初值和终值分别是多少？
    \item 计算线段 $(0,0) \to (10,10)$ 被窗口 $[2,5] \times [2,5]$ 裁剪后的端点。
    \item 写出完整的 Liang-Barsky 算法伪代码逻辑。
\end{enumerate}

\section{Bezier 曲线 (de Casteljau 算法)}
\begin{enumerate}
    \item 二阶 Bezier 曲线控制点为 $P_0(0,0), P_1(2,4), P_2(4,0)$。计算 $t=0.5$ 时的点。
    \item 三阶 Bezier 曲线控制点为 $P_0, P_1, P_2, P_3$。写出 $t=0.5$ 时的递推公式。
    \item 二阶 Bezier 曲线。计算 $t=0, t=1$ 时的点，并说明其物理意义。
    \item 证明 Bezier 曲线具有凸包性。
    \item 给定 $P_0(0,0), P_1(1,2), P_2(2,0)$，计算 $t=0.25$ 时的坐标。
    \item 什么是 Bernstein 基函数？写出三阶 Bezier 曲线的基函数表达式。
    \item Bezier 曲线的起点切线方向由哪些点决定？
    \item 绘制具有四个控制点的 Bezier 曲线的示意图，并标出特征多边形。
    \item 若将控制点 $P_1$ 平移，曲线会如何变化？（说明局部控制性问题）。
    \item 计算 $P_0(0,0), P_1(1,1), P_2(2,2), P_3(3,3)$ 定义的曲线（特殊情况）。
\end{enumerate}

\end{document}
